\documentclass[11pt, a4paper]{article}

\usepackage{amsmath}
\usepackage{graphicx}
\usepackage{booktabs}
\usepackage[hidelinks]{hyperref}
\usepackage{geometry}
\geometry{margin=2.5cm}

\title{Equalizer Lowpass FIR Filter Design via Convex Optimisation}
\author{}
\date{}

\begin{document}

\maketitle

\section{Overview}

A complex-valued lowpass FIR filter is designed whose passband target matches
the UCDC equalizer frequency response.  The filter is found by solving a
second-order cone programme (SOCP) using CVXPY.

\subsection*{Signal chain parameters}

\begin{tabular}{ll}
  \toprule
  Parameter & Value \\
  \midrule
  Sample rate $f_s$ & 750\,MHz (two-sided, Nyquist $= \pm375$\,MHz) \\
  Passband edge $f_p$ & $\pm162.5$\,MHz \\
  Stopband edge $f_{stop}$ & $\pm250.0$\,MHz \\
  Filter taps $N$ & 63 \\
  Passband tolerance $\delta_p$ & 0.05 \\
  Stopband tolerance $\delta_s$ & 0.1 \\
  \bottomrule
\end{tabular}


\section{Optimisation Formulation}

The DTFT of the complex tap vector $\mathbf{h} \in \mathbf{C}^N$ is
\[
  H(\omega) = \sum_{n=0}^{N-1} h[n]\, e^{-j\omega n}
  = \mathbf{a}(\omega)^\top \mathbf{h},
  \quad
  \mathbf{a}(\omega) = \bigl[1,\, e^{-j\omega},\, \ldots,\, e^{-j\omega(N-1)}\bigr]^\top.
\]

Sampling over a frequency grid turns this into the linear map
$\mathbf{H} = \mathbf{A}\mathbf{h}$, where
$A_{kn} = e^{-j\omega_k n}$.

Let $\mathcal{P}$, $\mathcal{S}$, and $\mathcal{T}$ denote the index sets of
the passband, stopband, and transition-band frequency grid points,
respectively.
The \textbf{fixed-tolerance} problem is
\begin{align}
  \underset{\mathbf{h}}{\text{minimise}} \quad
    & \sum_{k\in\mathcal{P}} \lvert \mathbf{a}(\omega_k)^\top \mathbf{h} - H_d(\omega_k) \rvert
      + \lambda \lVert \mathbf{h} \rVert^2 \notag \\
  \text{subject to} \quad
    & \lvert \mathbf{a}(\omega_k)^\top \mathbf{h} - H_d(\omega_k) \rvert \leq \delta_p,
      \quad k \in \mathcal{P} \notag \\
    & \lvert \mathbf{a}(\omega_k)^\top \mathbf{h} \rvert \leq \delta_s,
      \quad k \in \mathcal{S} \notag \\
    & \lvert \mathbf{a}(\omega_k)^\top \mathbf{h} \rvert \leq \delta_{t,k},
      \quad k \in \mathcal{T}. \notag
\end{align}

The transition-band bound $\delta_{t,k}$ decreases linearly from
$\max_{k\in\mathcal{P}}|H_d(\omega_k)|$ at the passband edge to
$3\delta_s$ at the stopband edge, preventing large overshoot in the
unconstrained gap while keeping the SOCP feasible.

\section{Filter Response}

Figure~\ref{fig:filter_response} shows the designed filter's magnitude
(dB) and phase response overlaid with the desired equalizer response
$H_d(\omega)$.

\begin{figure}[htbp]
  \centering
  \includegraphics[width=\linewidth]{../figure/eqz_lowpass/filter_response.png}
  \caption{Magnitude (top) and phase (bottom) of the designed equalizer
           lowpass FIR filter.  The green shaded regions indicate the
           passband $\pm162.5$\,MHz.  The dashed red line is the desired
           response $H_d(\omega)$ interpolated from the 256-bin UCDC
           equalizer FFT.}
  \label{fig:filter_response}
\end{figure}

Figure~\ref{fig:filter_response_inband} shows the in-band magnitude response
zoomed to $\pm5$\,MHz beyond the passband edges, with the y-axis limited to
$[-0.9, +0.1]$\,dB to resolve the passband ripple.

\begin{figure}[htbp]
  \centering
  \includegraphics[width=\linewidth]{../figure/eqz_lowpass/filter_response_inband.png}
  \caption{In-band magnitude response zoomed to
           $[\!-167.5, +167.5]$\,MHz with a $[-0.9, +0.1]$\,dB y-axis.
           The green shaded regions indicate the passband $\pm162.5$\,MHz.}
  \label{fig:filter_response_inband}
\end{figure}

\section{Complex Remez Algorithm}

The complex Remez algorithm (also known as Lawson's iteratively reweighted
least-squares, IRLS) extends the classical Parks--McClellan algorithm to
complex-valued FIR filters.  At each iteration the algorithm solves a
weighted least-squares problem
\[
  \min_{\mathbf{h}} \sum_k w_k(\ell)\,
    \bigl\lvert \mathbf{a}(\omega_k)^\top \mathbf{h} - H_d(\omega_k) \bigr\rvert^2,
\]
and then updates the weights as
$w_k(\ell+1) = w_k(\ell)\,\lvert e_k(\ell) \rvert$,
where $e_k(\ell)$ is the error at frequency $\omega_k$ after iteration
$\ell$.  This reweighting progressively concentrates the design effort on
the worst-case frequency points, driving the solution toward the
minimax (Chebyshev) optimum.

Unlike the convex optimisation approach, no explicit passband or stopband
tolerance parameters are required; the algorithm converges to the
equi-ripple solution determined solely by the frequency grid and the
desired response $H_d(\omega)$.

Figure~\ref{fig:complex_remez_response} shows the full-band magnitude and
phase response of the filter designed with the complex Remez algorithm.
Figure~\ref{fig:complex_remez_inband} zooms into the passband region, and
Figure~\ref{fig:complex_remez_error} shows the in-band approximation error.

\begin{figure}[htbp]
  \centering
  \includegraphics[width=\linewidth]{../figure/complex_remez_result/no_tb_constraint/filter_response.png}
  \caption{Full-band magnitude (top) and phase (bottom) of the complex
           Remez FIR filter.  The green shaded region indicates the
           passband $\pm162.5$\,MHz.  The dashed red line is the desired
           response $H_d(\omega)$.}
  \label{fig:complex_remez_response}
\end{figure}

\begin{figure}[htbp]
  \centering
  \includegraphics[width=\linewidth]{../figure/complex_remez_result/no_tb_constraint/filter_response_inband.png}
  \caption{In-band magnitude response of the complex Remez filter zoomed
           to the passband region.  The green shaded region indicates the
           passband $\pm162.5$\,MHz.}
  \label{fig:complex_remez_inband}
\end{figure}

\begin{figure}[htbp]
  \centering
  \includegraphics[width=\linewidth]{../figure/complex_remez_result/no_tb_constraint/inband_error.png}
  \caption{In-band approximation error of the complex Remez filter.
           The equi-ripple characteristic of the Chebyshev solution is
           visible across the passband.}
  \label{fig:complex_remez_error}
\end{figure}

\end{document}
